\section{Conclusion}\label{sec:conclusion}
% Summary of key findings
% Practical implications

This systematic review synthesized 9 studies published between 2022 and 2025 on the application of machine learning to the automated detection, classification, and feature extraction of Equatorial Plasma Bubbles. Three methodological paradigms were identified: end-to-end deep learning, hybrid dimensionality-reduction pipelines, and XAI-augmented frameworks. Reported accuracies ranged from 93\% to over 99\%, demonstrating the viability of ML as a replacement for manual observation. The review also revealed that the field remains concentrated on reactive detection using single-instrument optical data, with critical gaps in predictive spatiotemporal modeling and multimodal data fusion. Addressing the recurring limitations discussed herein---particularly environmental contamination, dataset imbalance, and limited cross-site generalizability---will be essential for transitioning these frameworks from research prototypes to operationally deployable systems.
